% Glossar
\chapter*{Glossar}
\vspace{-5em}
\addcontentsline{toc}{chapter}{Glossar}
\begin{itemize}
  \item \textbf{EcoHub-Plattform:} Eine zentrale Plattform zur digitalen Abwicklung von Geschäftsprozessen zwischen Versicherungen und externen Partnern wie Maklern.
  
  \item \textbf{Makler:} Eine externe Person oder Firma, die im Auftrag von Kunden Versicherungs- oder Vorsorgelösungen vermittelt.

  \item \textbf{Rollenzuweisung:} Der Prozess, bei dem Benutzern bestimmte Rechte oder Funktionen innerhalb einer Anwendung zugewiesen werden.

  \item \textbf{LDAP:} Ein Protokoll zur Abfrage und Modifikation von Verzeichnisdiensten über ein TCP/IP-Netzwerk.

  \item \textbf{API:} Eine definierte Schnittstelle, über die Anwendungen miteinander kommunizieren und Daten austauschen können.

  \item \textbf{Schnittstelle:} Eine definierte Verbindung zwischen zwei Systemen oder Komponenten zur Kommunikation und Datenübertragung.

  \item \textbf{Testabdeckung:} Ein Mass dafür, wie viel Prozent des Codes oder der Funktionen durch Tests geprüft werden.

  \item \textbf{CI/CD-Pipeline:} Eine automatisierte Prozesskette für das kontinuierliche Testen, Bauen und Deployen von Software.

  \item \textbf{Mapping-Funktion:} Eine Logik zur Umwandlung oder Zuordnung von Daten, z.\,B. zur Ermittlung der LDAP-Zuordnung basierend auf Benutzerrollen.

  \item \textbf{Applikationsgruppe:} Eine definierte Gruppe in einem Verzeichnisdienst, die Berechtigungen für bestimmte Anwendungen repräsentiert.

  \item \textbf{Playwright:} Ein Framework zur End-to-End-Testautomatisierung von Webanwendungen.

  \item \textbf{API-Endpunkt:} Eine konkrete Adresse innerhalb einer API, über die ein bestimmter Dienst aufgerufen werden kann.

  \item \textbf{Testautomation:} Die automatisierte Ausführung von Tests zur Überprüfung der Funktionalität und Qualität einer Software.

  \item \textbf{Heartbeat-Mechanismus:} Ein technischer Prüfmechanismus, der regelmässig die Erreichbarkeit eines Dienstes sicherstellt.

  \item \textbf{GitHub Actions:} Eine CI/CD-Plattform von GitHub zur Automatisierung von Softwareentwicklungs-Workflows.

  \item \textbf{YAML:} Eine menschenlesbare Auszeichnungssprache zur Konfiguration von Anwendungen und Systemen.

  \item \textbf{Rollenpräfix:} Ein vordefinierter Namensbestandteil, der zur Unterscheidung oder Kategorisierung von Benutzerrollen dient (z.\,B. \texttt{EL\_}, \texttt{KL\_}).

  \item \textbf{Einzelleben (EL):} Produktkategorie der Lebensversicherung für Einzelpersonen.

  \item \textbf{Kollektivleben (KL):} Produktkategorie der Lebensversicherung für Gruppen, z.\,B. Mitarbeitende eines Unternehmens.

  \item \textbf{Private Vorsorge (PV):} Freiwillige Vorsorge zur Ergänzung der staatlichen Absicherung.

  \item \textbf{Berufliche Vorsorge (BV):} Obligatorische Vorsorge für Arbeitnehmende zur Sicherung des gewohnten Lebensstandards im Alter.

  \item \textbf{Brokerportal:} Ein Webportal für Makler zur Abwicklung von Versicherungsprozessen.

  \item \textbf{Vertragservice:} Eine Funktionseinheit zur Verwaltung und Bearbeitung von Versicherungsverträgen.

  \item \textbf{Logismata B2B:} Eine B2B-Plattform oder Schnittstelle für die digitale Übermittlung von Versicherungsdaten.

  \item \textbf{QuickSale:} Eine Anwendung oder Funktion zur schnellen und einfachen Produktplatzierung oder Vertragsabwicklung.

  \item \textbf{Distinguished Name (DN):} Eindeutiger Bezeichner eines Objekts in einem LDAP-Verzeichnisbaum.

  \item \textbf{Payload:} Der Teil einer Nachricht, der die eigentlichen Daten transportiert, z.\,B. in einer API-Anfrage.

  \item \textbf{SOAP:} Ein XML-basiertes Protokoll zur Kommunikation zwischen Anwendungen über HTTP.

  \item \textbf{Single Sign-On (SSO):} Authentifizierungsverfahren, bei dem sich ein Nutzer einmal anmeldet und dann Zugriff auf mehrere Systeme hat.

  \item \textbf{Web Access Management (WAM):} Verwaltung von Zugriffskontrollen auf Webanwendungen und -dienste.

  \item \textbf{OpenLDAP:} Eine Open-Source-Implementierung des LDAP-Protokolls zur Verwaltung von Verzeichnisdiensten.

  \item \textbf{Docker Compose:} Ein Werkzeug zur Definition und Verwaltung von Multi-Container-Docker-Anwendungen über eine YAML-Datei.

  \item \textbf{End-to-End-Test (E2E):} Ein Test, der einen vollständigen Geschäftsprozess oder Anwendungsablauf von Anfang bis Ende prüft.

  \item \textbf{Testdaten:} Spezifisch vorbereitete Daten, die zur Ausführung von Tests verwendet werden.

  \item \textbf{Testumgebung:} Eine von der Produktivumgebung getrennte Infrastruktur zur Durchführung von Tests.

  \item \textbf{Testframework:} Ein Software-Werkzeug zur Strukturierung, Verwaltung und Ausführung automatisierter Tests.

  \item \textbf{3-2-1-1-0-Backup-Strategie:} Backup-Konzept mit 3 Kopien, 2 verschiedenen Medientypen, 1 externer Kopie, 1 Offline-Kopie und 0 Fehler bei der Wiederherstellung.

  \item \textbf{Backupkonzept:} Ein Plan zur Sicherung von Daten, um deren Wiederherstellung bei Verlust oder Ausfall zu gewährleisten.

  \item \textbf{Cloud-Speicherung:} Die Online-Speicherung von Daten auf externen Servern, die über das Internet erreichbar sind.

  \item \textbf{Commit:} Das Festschreiben von Änderungen in ein Versionskontrollsystem wie Git.

  \item \textbf{Continuous Deployment (CD):} Automatisiertes Ausrollen von geprüfter Software in die Produktionsumgebung.

  \item \textbf{Embedded LDAP:} Eine in eine Anwendung integrierte, lokale LDAP-Instanz zu Test- oder Entwicklungszwecken.

  \item \textbf{Git:} Ein verteiltes Versionskontrollsystem zur Verwaltung von Quellcode und Änderungen im Projektverlauf.

  \item \textbf{Google Drive:} Ein Cloud-Dienst von Google zur Speicherung und gemeinsamen Bearbeitung von Dateien.

  \item \textbf{Kanban:} Eine agile Projektmanagement-Methode zur Visualisierung und Steuerung von Arbeitsprozessen in kontinuierlichem Fluss.

  \item \textbf{Kanban-Board:} Ein visuelles Werkzeug zur Darstellung von Aufgaben und deren Bearbeitungsstatus (z.\,B. \textit{To Do}, \textit{In Progress}, \textit{Done}).

  \item \textbf{Meilenstein:} Ein definierter Punkt im Projektverlauf, an dem ein wesentliches Zwischenziel erreicht wird.

  \item \textbf{Organigramm:} Eine grafische Darstellung der Aufbauorganisation mit Rollen, Zuständigkeiten und Hierarchien.

  \item \textbf{Projektaufbauorganisation:} Die strukturierte Darstellung von Zuständigkeiten und Rollen innerhalb eines Projekts.

  \item \textbf{Projektphase:} Eine zeitlich abgegrenzte Etappe im Projekt, z.\,B. Analyse, Planung/Design, Implementierung, Test/Qualitätssicherung oder Betrieb/Wartung.

  \item \textbf{Projektstrukturplan:} Eine hierarchische Zerlegung eines Projekts in Teilaufgaben und Arbeitspakete zur besseren Planbarkeit.

  \item \textbf{Scrum:} Ein agiles Framework für Projektmanagement und Softwareentwicklung mit iterativen Sprints und klar definierten Rollen.

  \item \textbf{Single Point of Failure:} Eine Komponente, deren Ausfall das gesamte System oder einen kritischen Dienst lahmlegen kann.

  \item \textbf{Sprint:} Eine zeitlich begrenzte Arbeitsperiode im Scrum-Framework, in der ein definiertes Ziel erreicht werden soll.

  \item \textbf{Stakeholder:} Alle Personen oder Gruppen, die direkt oder indirekt vom Projekt betroffen sind oder Einfluss darauf haben.

  \item \textbf{Testkonzept:} Ein übergeordnetes Dokument, das die Ziele, Strategien, Methoden und Zuständigkeiten für das Testen festlegt.

  \item \textbf{Teststrategie:} Ein strukturierter Plan zur Auswahl und Durchführung geeigneter Testarten zur Qualitätssicherung.

  \item \textbf{Unit-Test, Component-Test, Integration-Test:} Unterschiedliche Testarten zur Überprüfung einzelner Funktionen (Unit), Komponenten (Component) und deren Zusammenspiel (Integration).

  \item \textbf{Versionierung, Versionskontrolle:} Die systematische Nachverfolgung und Verwaltung von Änderungen an Dateien oder Quellcode über Zeit hinweg.

  \item \textbf{Wasserfallmodell:} Ein lineares Projektvorgehensmodell mit aufeinanderfolgenden Phasen ohne Rücksprünge.

\end{itemize} 